\section{Resultados documentais}
\subsection{Testemunhos confirmados e suas datas}
A Tabela~\ref{tab:attestations} reúne passagens selecionadas com localizadores. Ela não é uma série de frequências e não representa a proporção de brasileiros que empregava cada grafia. As datas literárias identificam a obra; quando o acesso é uma transcrição posterior, isso é explicitado. As datas de estudos acadêmicos não são convertidas automaticamente em datas das interações.

\small
\begin{longtable}{p{0.07\textwidth}p{0.24\textwidth}p{0.17\textwidth}p{0.37\textwidth}}
\caption{Atestações selecionadas: ano de referência, grafia e localizador. E1: imagem; E2: transcrição primária; E3: interação reportada.}\label{tab:attestations}\\
\toprule Ano & Documento & Grafia & Localizador / evidência \\\midrule\endfirsthead
\toprule Ano & Documento & Grafia & Localizador / evidência \\\midrule\endhead
1872 & Til \cite{alencar1872} & Quiá!... quiá!... quiá!... & Volume II, capítulo VIII; âncora: Ai, menina! (E2). \\
1879 & A Imprensa Ytuana : jornal scientifico, litterario, noticiozo e industrial do Instituto do Novo-Mundo, ano 004, n. 183, 1879 \cite{r3doc1} & (Riso) & 1879-09-13; PDF p. 2 (E1). \\
1884 & A segunda vida \cite{assis1884} & cá, cá, cá, cá, cá... & Âncora: Era uma gargalhada assim (E2). \\
1900 & Anais do Senado Federal, livro 1 \cite{anual1900} & (Riso.) & PDF p. 248 (E2). \\
1901 & A cidade e as serras \cite{queiros1901} & Ah! ah! ah! & Âncora: Então tambem eu rio!; PG 18220 (E2). \\
1902 & Anais do Senado Federal, livro 1 \cite{anual1902} & (Riso) & PDF p. 60 (E2). \\
1903 & Anais do Senado Federal, livro 1 \cite{anual1903} & (Riso.) & PDF p. 114 (E2). \\
1904 & Anais do Senado Federal, livro 1 \cite{anual1904} & (Riso.) & PDF p. 320 (E2). \\
1905 & Anais do Senado Federal, livro 1 \cite{anual1905} & (riso) & PDF p. 88 (E2). \\
1906 & Anais do Senado Federal, livro 1 \cite{anual1906} & (Risos) & PDF p. 318 (E2). \\
1907 & Anais do Senado Federal, livro 1 \cite{anual1907} & (Riso.) & PDF p. 225 (E2). \\
1908 & Anais do Senado Federal, livro 3 \cite{anual1908} & (Riso.) & PDF p. 169 (E2). \\
1909 & Anais do Senado Federal, livro 2 \cite{anual1909} & (Riso.) & PDF p. 21 (E2). \\
1910 & Anais do Senado Federal, livro 2 \cite{anual1910} & (Riso) & PDF p. 153 (E2). \\
1911 & Anais do Senado Federal, livro 1 \cite{anual1911} & (Riso nas galerias.) & PDF p. 312 (E2). \\
1912 & Anais do Senado Federal, livro 1 \cite{anual1912} & (Risos.) & PDF p. 418 (E2). \\
1913 & Anais do Senado Federal, livro 5 \cite{anual1913} & (Risos nas galerias.) & p. 238, PDF p. 255 (E2). \\
1914 & Anais do Senado Federal, livro 1 \cite{anual1914} & (Risos.) & PDF p. 205 (E2). \\
1915 & Anais do Senado Federal, livro 1 \cite{anual1915} & (riso) & PDF p. 121 (E2). \\
1916 & Anais do Senado Federal, livro 4 \cite{anual1916} & (Risos.) & sessão de 24/08/1916, p. 399, PDF p. 415 (E2). \\
1917 & A engomadeira \cite{almada1917} & Ah! Ah! Ah!... & Âncora: E 'inda o confessa!; PG 23879 (E2). \\
1918 & Diário da Câmara dos Deputados, n. 008 \cite{anual1918} & Risos. & data no catálogo 1918-08-02, p. 18 (E2). \\
1919 & Urupês, quinta edição \cite{lobato1919} & Quá, quá, quá! & p. impressa 19; PDF p. 24; conto O engraçado arrependido (E1). \\
1920 & Diário da Câmara dos Deputados, n. 059 \cite{anual1920} & Risos. & data no catálogo 1920-03-31, p. 41 (E2). \\
1921 & Anais do Senado Federal, livro 1 \cite{anual1921} & (Riso.) & PDF p. 500 (E2). \\
1922 & Anais do Senado Federal, livro 3 \cite{anual1922} & (Riso.) & PDF p. 12 (E2). \\
1923 & Anais do Senado Federal, livro 2 \cite{anual1923} & (Riso.) & PDF p. 92 (E2). \\
1924 & Anais do Senado Federal, livro 1 \cite{anual1924} & (Riso.) & PDF p. 611 (E2). \\
1925 & Diário da Câmara dos Deputados, n. 77 \cite{portugal1925} & Risos. & 11/06/1925; p. 14; intervenção de Cunha Leal (E2). \\
1926 & Anais do Senado Federal, livro 1 \cite{anual1926} & (Risos.) & PDF p. 311 (E2). \\
1927 & Fon-Fon \cite{anual1927} & Quá, quá & PDF p. 75; Saibam todos (E1). \\
1928 & Diário Nacional, n. 290 \cite{diarionacional1928} & (risos, muito bem) & 17/06/1928; p. 4; discurso de Mattos Pimenta (E1). \\
1929 & O Tico-Tico \cite{anual1929} & Quá, quá, quá, quá.... & 12/06/1929; PDF p. 16; A raposa e o gallo; riso ficcional explícito (E1). \\
1930 & A Batalha \cite{anual1930} & Quá-quá-quá...! & 06/08/1930, PDF p. 2, seção A pedido (E1). \\
1931 & O Tico-Tico \cite{anual1931} & quá quá, quá... & 22/04/1931; PDF p. 23; Ri... quá quá, quá... (E1). \\
1932 & Diario de Noticias \cite{anual1932} & Quá-quá! & 02/10/1932; PDF p. 12; comentário de xadrez assinado Pharaó (E1). \\
1933 & Anais da Assembleia Nacional Constituinte, livro 1 \cite{anual1933} & (Riso.) & PDF p. 318 (E2). \\
1934 & Anais da Assembleia Nacional Constituinte, livro 6 \cite{anual1934} & (Risos.) & PDF p. 79 (E2). \\
1935 & O Tico-Tico, n. 1569 \cite{ticotico1935} & Quá quá quá... & 30/10/1935; capa, último quadro; PDF p. 1 (E1). \\
1936 & Diário das Sessões da Assembleia Nacional, n. 075 \cite{anual1936} & (Risos) & sessão 1936-02-07; p. 437 (E2). \\
1937 & Diário das Sessões da Assembleia Nacional, n. 117 \cite{anual1937} & (Risos) & sessão 1937-02-11; p. 443 (E2). \\
1938 & O Tico-Tico \cite{anual1938} & Ha ha ha ha! & 01/06/1938; PDF p. 18; As aventuras do Camondongo Mickey (E1). \\
1939 & Diário das Sessões da Assembleia Nacional, n. 053 \cite{anual1939} & (Risos) & sessão 1939-12-11; p. 33 (E2). \\
1940 & O Tico-Tico \cite{anual1940} & Quá - quá! quá! & 22/05/1940; PDF p. 26; As aventuras do Camondongo Mickey (E1). \\
1941 & Diário Carioca, n. 3902 \cite{carioca1941} & Ah! ah! & 08/03/1941; p. 13; folhetim As memórias de Sherlock Holmes (E1). \\
1942 & Diário das Sessões da Assembleia Nacional, n. 008 \cite{anual1942} & Risos. & data no catálogo 1942-12-12, p. 49 (E2). \\
1943 & Diário das Sessões da Assembleia Nacional, n. 024 \cite{anual1943} & Risos. & data no catálogo 1943-03-17, p. 235 (E2). \\
1944 & Diário das Sessões da Assembleia Nacional, n. 066 \cite{anual1944} & (Risos) & sessão 1944-03-22; p. 298 (E2). \\
1945 & Diário das Sessões da Assembleia Nacional, n. 118 \cite{anual1945} & (risos) & sessão 1945-01-11; p. 115 (E2). \\
1946 & Anais do Senado Federal, livro 1 \cite{anual1946} & (Risos.) & PDF p. 76 (E2). \\
1947 & Jornal de Notícias, n. 267 \cite{jornalnoticias1947} & (Risos dos presentes) & 02/03/1947; p. 1; fala de Bidault; não é Folha da Manhã (E1). \\
1948 & Anais do Senado Federal, livro 2 \cite{anual1948} & (Risos) & PDF p. 501 (E2). \\
1949 & Diário das Sessões da Assembleia Nacional, n. 180 \cite{anual1949} & Risos. & sessão 1949-03-24; p. 325 (E2). \\
1950 & Carioca, n. 792 \cite{anual1950} & Quá, quá, quá & 07/12/1950, PDF p. 48, coluna Por trás do dial (E1). \\
1951 & Diário das Sessões da Assembleia Nacional, n. 105 \cite{anual1951} & Risos. & data no catálogo 1951-04-28, p. 984 (E2). \\
1952 & Diário das Sessões da Assembleia Nacional, n. 132 \cite{anual1952} & (risos) & data no catálogo 1952-03-04, p. 427 (E2). \\
1953 & Anais do Senado Federal, livro 1 \cite{anual1953} & (Risos). & PDF p. 14 (E2). \\
1954 & Anais do Senado Federal, livro 1 \cite{anual1954} & (Riso). & PDF p. 7 (E2). \\
1955 & Diário do Congresso Nacional, suplemento \cite{camara1955} & (Risos) & 10/12/1955; p. 13; intervenção de Alberto Torres (E1). \\
1956 & Mundo Esportivo \cite{anual1956} & quá, quá, quá & 25/05/1956, PDF p. 2, diálogo humorístico (E1). \\
1957 & Anais do Senado Federal, livro 1 \cite{anual1957} & (Risos). & PDF p. 230 (E2). \\
1958 & Anais do Senado Federal, livro 1 \cite{anual1958} & (Riso). & PDF p. 662 (E2). \\
1959 & Anais do Senado Federal, livro 2 \cite{anual1959} & (Riso) & PDF p. 398 (E2). \\
1960 & Anais do Senado Federal, livro 2 \cite{anual1960} & (Riso) & PDF p. 219 (E2). \\
1961 & Anais do Senado Federal, livro 1 \cite{anual1961} & (Riso). & PDF p. 159 (E2). \\
1962 & Anais do Senado Federal, livro 1 \cite{anual1962} & (Riso.) & PDF p. 190 (E2). \\
1963 & Diário das Sessões da Assembleia Nacional, n. 066 \cite{anual1963} & (Risos) & data no catálogo 1963-01-08, p. 1747 (E2). \\
1964 & Anais do Senado Federal, livro 2 \cite{anual1964} & (Riso) & PDF p. 25 (E2). \\
1965 & Anais do Senado, livro 1 \cite{senado1965} & (Risos.) & p. impressa 162; PDF p. 185; intervenção de Edmundo Levi (E1). \\
1966 & Anais do Senado Federal, livro 1 \cite{anual1966} & (Riso) & PDF p. 148 (E2). \\
1967 & Diário das Sessões da Assembleia Nacional, n. 059 \cite{anual1967} & (Risos) & sessão 1967-01-19; p. 1059 (E2). \\
1968 & Anais do Senado Federal, livro 1 \cite{anual1968} & (Risos.) & PDF p. 137 (E2). \\
1969 & Anais do Senado Federal, livro 1 \cite{anual1969} & (Riso.) & PDF p. 164 (E2). \\
1970 & Anais do Senado Federal, livro 4 \cite{anual1970} & (Risos.) & p. 215, PDF p. 216 (E2). \\
1971 & Anais do Senado Federal, livro 1 \cite{anual1971} & (Risos.) & PDF p. 233 (E2). \\
1972 & Anais do Senado Federal, livro 1 \cite{anual1972} & (risos) & PDF p. 68 (E2). \\
1973 & Anais do Senado Federal, livro 1 \cite{anual1973} & (Risos.) & PDF p. 124 (E2). \\
1974 & Anais do Senado Federal, livro 1 \cite{anual1974} & (Risos.) & PDF p. 344 (E2). \\
1975 & Anais do Senado Federal, livro 1 \cite{anual1975} & (risos) & PDF p. 173 (E2). \\
1976 & Anais do Senado Federal, livro 1 \cite{anual1976} & (Risos.) & PDF p. 49 (E2). \\
1977 & Anais do Senado Federal, livro 1 \cite{anual1977} & (Risos.) & PDF p. 83 (E2). \\
1978 & Anais do Senado Federal, livro 1 \cite{anual1978} & (Risos.) & PDF p. 56 (E2). \\
1979 & Correio do Povo \cite{correio1979} & AH, AH, AH & 28/04/1979; p. 5; tira Bidu, quadros finais (E1). \\
1980 & Anais do Senado Federal, livro 1 \cite{anual1980} & (Risos.) & PDF p. 35 (E2). \\
1981 & Anais do Senado Federal, livro 2 \cite{anual1981} & (risos) & PDF p. 81 (E2). \\
1982 & Anais do Senado Federal, livro 1 \cite{anual1982} & (Risos.) & PDF p. 100 (E2). \\
1983 & Anais do Senado Federal, livro 1 \cite{anual1983} & (Risos.) & PDF p. 57 (E2). \\
1984 & Anais do Senado Federal, livro 3 \cite{anual1984} & (Risos.) & PDF p. 281 (E2). \\
1985 & Diário do Congresso Nacional, seção I \cite{camara1985} & (risos e aplausos.) & 29/03/1985; p. impressa 2107; PDF p. 55 (E1). \\
1986 & Anais do Senado Federal, livro 3 \cite{anual1986} & (Risos.) & PDF p. 13 (E2). \\
1987 & Anais do Senado Federal, livro 1 \cite{anual1987} & (Risos.) & PDF p. 129 (E2). \\
1988 & Anais do Senado Federal, livro 1 \cite{anual1988} & (risos.) & PDF p. 5 (E2). \\
1989 & Anais do Senado Federal, livro 1 \cite{anual1989} & (Risos) & PDF p. 82 (E2). \\
1990 & Anais do Senado Federal, livro 1 \cite{anual1990} & (Risos) & PDF p. 48 (E2). \\
1991 & Anais do Senado Federal, livro 1 \cite{anual1991} & (Risos.) & PDF p. 43 (E2). \\
1992 & Anais do Senado Federal, livro 2 \cite{anual1992} & (Risos) & diário de 20/02/1992, p. 358, PDF p. 110 (E2). \\
1993 & Diário da Assembleia da República, I série, n. 073 \cite{anual1993} & Risos do PSD e do CDS. & data no catálogo 1993-05-20, p. 2300 (E2). \\
1994 & Folha de S.Paulo \cite{folha1994} & (risos) & 03/09/1994; Veja íntegra da conversa de Ricupero com Monforte (E2). \\
1995 & Anais do Senado Federal, livro 1 \cite{anual1995} & (Risos) & PDF p. 144 (E2). \\
1996 & Diário da Assembleia da República, I série, n. 032 \cite{anual1996} & Risos. & data no catálogo 1996-01-26, p. 911 (E2). \\
1997 & Anais do Senado Federal, livro 2 \cite{anual1997} & (risos) & janeiro de 1997, entrevista reproduzida nos Anais, p. 13, PDF p. 39 (E2). \\
1998 & Diário da Assembleia da República, I série, n. 002 \cite{anual1998} & Risos do PSD. & data no catálogo 1998-09-17, p. 40 (E2). \\
1999 & Anais do Senado Federal, livro 1 \cite{anual1999} & (Risos) & PDF p. 507 (E2). \\
2000 & Diário da Assembleia da República, I série, n. 033 \cite{anual2000} & Risos do PCP. & data no catálogo 2000-12-21, p. 1292 (E2). \\
2001 & Diário da Assembleia da República, I série, n. 034 \cite{anual2001} & Risos do CDS-PP e do PSD. & data no catálogo 2001-01-03, p. 1372 (E2). \\
2002 & Diário da Assembleia da República, I série, n. 031 \cite{anual2002} & Risos. & data no catálogo 2002-01-23, p. 1205 (E2). \\
2003 & Diário da Assembleia da República, I série, n. 071 \cite{anual2003} & Risos do PCP. & data no catálogo 2003-01-08, p. 3008 (E2). \\
2004 & Diário da Assembleia da República, I série, n. 035 \cite{anual2004} & Risos do PCP. & data no catálogo 2004-01-07, p. 2001 (E2). \\
2005 & Anais do Senado Federal, livro 1 \cite{anual2005} & (Risos.) & PDF p. 606 (E2). \\
2006 & Anais do Senado Federal, livro 1 \cite{anual2006} & (Risos.) & PDF p. 361 (E2). \\
2007 & Chat reproduzido por Young (2008) \cite{young2008} & kkkk & Dissertação p. 93; contexto: semestre 2007.2 (E3). \\
2008 & Anais do Senado Federal, livro 1 \cite{anual2008} & (Risos.) & PDF p. 957 (E2). \\
2009 & Diário da Assembleia da República, I série, n. 030 \cite{anual2009} & Risos do PS. & data no catálogo 2009-01-07, p. 49 (E2). \\
2010 & Leão (2010) \cite{leao2010} & kkkkkkkk / rsrsrsrs & PDF pp. 90-92; fragmentos 97, 101 e 103 e análise (E3). \\
2010 & Cifra Club: como afinar meio tom abaixo \cite{r4forumcifraclub} & kkkkkkkkkkkkk & resposta Enviado: 12/jul/10 15:21; após boa!; captura 2010-09-15T13:54:14Z (E2). \\
2014 & Orkut: comunidade Kiko Dança Psy-TRANCE \cite{r4orkutcomunidades} & hehehe & descrição da comunidade, após sempre agitado; captura 2014-10-04T09:58:23Z (E2). \\
2015 & Silva (2015) \cite{silva2015} & Kkkk & p. impressa 86, PDF p. 87; CA7 P2 (E3). \\
2016 & A Cidade de Jardinópolis \cite{r3doc18} & Hahaha... & 2016-01-09; pasta 49, página lógica 3; p. impressa 3; Só Rindo, Claudio Aguilhera (E1). \\
2019 & Macunaíma, edição UFFS \cite{andrade2019} & Quá! quá! quá quaquá!... & p. impressa e PDF 139; ficha catalográfica PDF 205 (E1). \\
2022 & A Cidade de Jardinópolis \cite{r3doc19} & kkkkk & 2022-01-29; pasta 57, página lógica 21; p. impressa 5; Entrevista (E1). \\
2025 & Correio Braziliense: Uma sintonia capital \cite{r3doc21} & (risos) & 2025-04-07; PDF p.22; Diversão \& Arte; entrevista de Walter Lima Jr. por Ricardo Daehn (E1). \\
\bottomrule\end{longtable}\normalsize

Em \textit{Til}, segundo volume, capítulo VIII, a reação de Nhá Tudinha contém ``Quiá!... quiá!... quiá!...'' \cite{alencar1872}. O contexto descreve uma risada e permite classificar a sequência como representação de riso. A passagem confirma a grafia no texto consultado, mas não estabelece que seja a primeira onomatopeia de riso em português. Também não contém \forma{kkkk}. A autoria é de José de Alencar; atribuí-la genericamente a Eça é um erro de identificação.

Em \textit{A segunda vida}, de Machado de Assis, o corpus eletrônico consultado apresenta ``cá, cá, cá, cá, cá...'' após a descrição de uma gargalhada \cite{assis1884}. O caso mostra a relação entre encenação literária e representação sonora. A presença do mesmo mecanismo repetitivo não revela por si só se o autor conhecia determinada passagem de outro escritor ou se um usuário de chat a herdou mais tarde.

Em \textit{A cidade e as serras}, publicado em 1901, encontra-se uma sequência de \forma{ah} associada explicitamente a rir \cite{queiros1901}. A edição eletrônica permite localizar a passagem pela âncora ``Então tambem eu rio!'' ou sua variante efetiva indicada no inventário. Em \textit{A engomadeira}, de 1917, a sequência ``Ah! Ah! Ah!...'' integra uma cena de zombaria \cite{almada1917}. Esses exemplos impedem reduzir a história do riso a uma sucessão em que um formato substitui integralmente outro a cada década.

Na capa de \textit{O Tico-Tico}, número 1569, de 30 de outubro de 1935, a narrativa gráfica ``A caricatura do Bolinha'' apresenta \forma{quá quá quá} no último quadro \cite{ticotico1935}. A imagem foi inspecionada para resolver ruídos do OCR. O testemunho é particularmente útil porque reúne data, número, localização e contexto visual. Demonstra uso numa revista infantil brasileira antes da cronologia que atribuía sua consolidação inicial aos quadrinhos Disney dos anos 1950 ou 1960; não determina, isoladamente, o alcance social da forma.


O \textit{Diário Carioca}, em 8 de março de 1941, publica uma sequência de \forma{ah} no folhetim traduzido de Conan Doyle, página 13 \cite{carioca1941}. A data identifica a publicação brasileira da tradução, não a composição original. No \textit{Correio do Povo} de 28 de abril de 1979, página 5, a tira do Bidu contém repetições de \forma{AH} nos quadros finais \cite{correio1979}. A imagem permite corrigir o OCR, que alternava equivocadamente \forma{AU} e \forma{AH}. O caso demonstra permanência dessa convenção em quadrinhos de 1979, sem medir frequência no conjunto da imprensa.

\subsection{A documentação digital de 2007 e os estudos posteriores}
\subsubsection{Rubricas institucionais e jornalísticas: 1925 a 1994}
As rubricas foram incluídas em categoria própria porque representam uma reação descrita, não uma imitação fonética. No Diário da Câmara dos Deputados de Portugal, sessão de 11 de junho de 1925, página 14, a intervenção de Cunha Leal é seguida de \forma{Risos.} na transcrição oficial \cite{portugal1925}. O portal informa que a extração é automática; o testemunho foi classificado como E2, sem equiparação a colação do fac-símile.

No suplemento do Diário do Congresso Nacional de 10 de dezembro de 1955, página 13, a intervenção de Alberto Torres recebe a marca \forma{(Risos)} \cite{camara1955}. Nos Anais do Senado de 1965, livro 1, página impressa 162 e página 185 do PDF, a fala de Edmundo Levi sobre a liberdade dos senadores e um cafezinho é acompanhada de \forma{(Risos.)} \cite{senado1965}. A imagem confirma o sinal; o dia da sessão não foi individualizado neste registro, preservando-se apenas o ano do volume.

No Diário do Congresso Nacional de 29 de março de 1985, página 2107, a marca \forma{(risos e aplausos.)} aparece em material transcrito inserido nos trabalhos \cite{camara1985}. Na versão eletrônica da Folha de S.Paulo de 3 de setembro de 1994, a transcrição da conversa de Ricupero com Monforte inclui \forma{(risos)} após uma anedota \cite{folha1994}. Esses dados cobrem modalidades descritivas em diferentes décadas; não autorizam calcular predominância nem equiparar rubrica e onomatopeia.

Com as novas atestações, o inventário tem ao menos um documento de cada década de 1900 a 2000. Essa cobertura decenal mínima não significa cobertura de todos os anos, três livros e três periódicos por período, nem amostra equilibrada entre suportes. As tabelas anuais explicitam a diferença.

\subsubsection{Interações digitais reproduzidas em pesquisas}
Young \cite{young2008} reproduz, na página 93, um diálogo de bate-papo com cadeias de \forma{k}; a pesquisa é situada no semestre 2007.2. O registro é um testemunho reportado de interação digital, acessível em fonte acadêmica de 2008. O dado autoriza afirmar que \forma{kkkk} é documentado naquele conjunto ligado a 2007. Não autoriza concluir que surgiu em 2007 ou que aquela plataforma tenha sido seu berço nacional. O download direto do PDF falhou no ambiente local, embora o documento e o registro institucional tenham sido consultados pelo serviço de leitura web; essa diferença consta no inventário.

Leão \cite{leao2010} reúne, em pesquisa publicada em 2010, exemplos de chat com \forma{kkkk}, \forma{rsrsrsrs} e sequências heterogêneas, nas páginas de arquivo 90--92. Silva \cite{silva2015} registra \forma{Kkkk} no comentário CA7 P2, página impressa 86, correspondente à página 87 do PDF, e documenta outras formas nos anexos. Esses casos mostram coexistência de repertórios em contextos educacionais. Não constituem prova de que um tipo tenha desaparecido nacionalmente quando outro passou a circular.

O ano de 2010 ou 2015 é o ano do documento consultado quando não há data individual validada da mensagem. Preservar essa distinção evita construir uma linha temporal com precisão inexistente. Do mesmo modo, um artigo recente que reproduz comentários antigos deve fornecer metadados da coleta para que os comentários sejam datados independentemente da publicação.

\subsection{Busca, downloads e falsos positivos}
O conjunto baixado inclui textos literários sem sequência de riso confirmada pela busca dirigida, PDFs de periódicos e documentos de apoio. Entre os resultados de busca, ``Quá... Quá...'' aparece também como nome comercial em \textit{Jornal de Notícias}, de 1946, e como parte do título de um programa na revista \textit{Intervalo}, de 1965. Esses usos podem ser pertinentes à circulação cultural de uma expressão, mas não são contados como transcrições de alguém rindo.

As buscas em três números de \textit{O Malho}, de 1920, não produziram atestação que tivesse sido validada para inclusão na tabela principal. Uma correspondência automática atravessava a palavra \forma{ahi} e a forma verbal \forma{ha}, gerando um falso positivo. O resultado demonstra a necessidade de fronteiras de palavra e conferência contextual. Não demonstra ausência de riso onomatopaico em todo o periódico ou na década.

\subsection{A auditoria da série de 1900 a 2000}
O material inicial apresenta onze blocos decenais, cada um com exemplo literário e jornalístico supostamente literal. Sua apresentação dá aparência de uma série documental contínua, mas faltam páginas, edições específicas e, em vários casos, identificação suficiente de autor, seção ou número. A matriz no apêndice conserva os 22 alvos de conferência. Nenhuma dessas frases é promovida a citação autêntica apenas pela existência da obra ou do periódico citado.

O problema não é apenas bibliográfico. A série associa 1900 a solenidade, 1910 a caricatura, 1920 a riso de salão, 1930 a escárnio, 1940 a rubricas, 1950 a Disney, 1960 a contracultura, 1970 a gibis, 1980 a datilografia, 1990 a \forma{rs} e 2000 a \forma{kkkk}. Sem contagens comparáveis, trata-se de uma narrativa interpretativa, não de uma distribuição observada. A existência de rubricas institucionais ou de quadrinhos numa época não implica predomínio nacional de uma grafia.

As alegações de que a norma exigia sempre \forma{h} final, de que ninguém agrupava letras em 1900 e de que o \forma{rs} nasceu para economizar fita datilográfica não receberam documentação normativa ou material suficiente. A explicação econômica é testável, mas não constitui prova em si mesma. Os exemplos inventados para ilustrar um estilo, quando apresentados como transcrição, devem ser removidos da condição de fonte histórica.
